\documentclass[../main.tex]{subfiles}
\graphicspath{{\subfix{../images/}}}

\begin{document}

\section{Introduction}
As climate change has accelerated at an unprecedented rate,\citep{nasa_evidence_2024} it has driven rising temperatures and increased variability in precipitation patterns.\citep{westerling2006warming} Higher temperatures and prolonged droughts present a unique risk for wildfire activity. Simultaneously, a substantial increase in the number of people living high wildfire risk areas, such as the wildland-urban interface (WUI), has quickly elevated the number of people at risk.\citep{burke2021changing, marlon2012long} This has led to increased fire suppression efforts, yielding a substantial increase in available wildfire fuel.\citep{burke2021changing, marlon2012long} Together, these have increased wildfire frequency and severity in the western United States.\citep{casey2024measuring, hurteau2014projected, dennison2014large, radeloff2018rapid, abatzoglou2016impact, burke2021changing} 

Wildfires have profound impacts on human health and communities.\citep{casey2024measuring, moritz2014learning} Wildfires and the \PM from wildfire smoke are associated with adverse cardiorespiratory, neurological, and mental health outcomes.\citep{mcbrien_wildfire_2023, elser2025wildfire, casey2024measuring, epa_pm_2019} Community-level effects of wildfires include displacement, power outages, and damage and destruction to communities.\citep{mcbrien_wildfire_2023, gill2013worldwide}

Of these community impacts, the risk of wildfire and power outage co-occurrence presents a unique challenge. Wildfires may cause power outages, and power outages may cause wildfires. Wildfires cause power outages via high winds that cause falling tree branches on lines or that cause electric poles or lines to fall onto dry ground\citep{khan2022probabilistic}. Wildfires may also cause excessive heat that can weaken power lines.\citep{shi2018model} Power outages may cause wildfires when electrical infrastructure (e.g., power lines) comes into contact with trees and branches -- typically as a result of either the power line or the tree falling.\citep{sayarshad2023preignition} 

As wildfire frequency continues to increase, the co-occurrence of wildfires and power outages is increasing.\citep{do2023spatiotemporal, do2025spatiotemporal, mukherjee2018multi} Preliminary studies of the co-occurrence of climate hazards show that the risks associated with environmental hazards, including of wildfires and power outages, increase substantially when there are two or more simultaneous hazards.\citep{do2025spatiotemporal, mukherjee2018multi} This necessitates adoption of adaptive measures to protect public safety.

One such measure, implemented in California and increasingly adopted by other states, are Public Safety Power Shutoffs (PSPS).\citep{psps_2025} During a PSPS event, electric utility companies proactively de-energize power lines to prevent power lines from causing or worsening wildfires.\citep{psps_2025} Intended as a preventative measure to protect people, PSPS events bring many of the same adverse health effects as power outages.\citep{psps_2025} Studies show that power outages are associated with adverse health outcomes, including injuries, hypothermia, hyperthermia, and mortality.\citep{do2025spatiotemporal} They can increase adverse health outcomes by exposing individuals to more extreme temperatures in the absence of air conditioning and heating systems or by preventing individuals from using air filtration systems during events that compromise air quality.\citep{do2025spatiotemporal} 

While the intention is to increase public safety, PSPS events result in power outages that may harm the same people they attempt to protect from wildfires. PSPS events may pose a particular threat when wildfires still occur despite the shutoff, exposing individuals to a dual hazard.\citep{do2025spatiotemporal} This may happen if (1) a PSPS event is used to prevent a wildfire from getting worse, or (2) a PSPS event is initiated to attempt to reduce already high wildfire risk but the wildfire still occurs. In both scenarios, PSPS events may result in dual exposure to environmental hazards. Studies of co-occurring hazards show that risk rises substantially when two or more hazards co-occur.\citep{do2023spatiotemporal, mukherjee2018multi} Wildfires co-occurring with power outages, in particular, present an elevated risk due to the high concentration of \PM in the air, the high stress levels associated with wildfires -- coupled with the inability to use air filtration systems or other electricity-dependent devices.

Despite the known health risks of wildfires, power outages, and wildfires co-occurring with power outages, few studies have examined the effect of PSPS events on health. This study aimed to describe the extent of dual wildfire-PSPS exposure and the association between this dual exposure and adverse health outcomes.

\bigskip
\bigskip
\end{document}